% Part: sets-functions-relations
% Chapter: relations
% Section: equivalence-relations

\documentclass[../../../include/open-logic-section]{subfiles}

\begin{document}

\olfileid{sfr}{rel}{eqv}

\olsection{Equivalentierelaties}

De identiteitsrelatie op een verzameling is reflexief, symmetrisch en
transitief. Relaties~$R$ met deze drie eigenschappen komen veel voor.

\begin{defn}[Equivalentierelatie] 
Een relatie $R \subseteq A^2$ die reflexief, symmetrisch en
transitief is, heet een \emph{equivalentierelatie}. !!^{element}s $x$
en $y$ uit~$A$ heten \emph{$R$-equivalent} als~$Rxy$.
\end{defn}

Equivalentierelaties leiden tot het begrip \emph{equivalentieklasse}.
Een equivalentierelatie deelt haar domein op in afzonderlijke delen.
Binnen ieder deel zijn alle objecten onderling gerelateerd; objecten
uit verschillende delen zijn niet aan elkaar gerelateerd. Soms is het
handig \emph{rechtstreeks} over deze delen te spreken. Daarom voeren
we de volgende definitie in:

\begin{defn}\ollabel{def:equivalenceclass}
Laat $R \subseteq A^2$ een equivalentierelatie zijn. Voor iedere $x \in A$
is de \emph{equivalentieklasse} van $x$ in~$A$ de verzameling $\equivrep{x}{R}
= \Setabs{y \in A}{Rxy}$. De \emph{quotiëntverzameling} van $A$ modulo~$R$ is
$\equivclass{A}{R} = \Setabs{\equivrep{x}{R}}{x \in A}$, dat wil zeggen de
verzameling van deze equivalentieklassen. 
\end{defn}

De volgende propositie rechtvaardigt de definitie van een
equivalentieklasse: zij laat zien dat de equivalentieklassen inderdaad
een partitie van~$A$ vormen:

\begin{prop}
Als $R \subseteq A^2$ een equivalentierelatie is, dan geldt $Rxy$ precies wanneer
$\equivrep{x}{R} = \equivrep{y}{R}$.
\end{prop}

\begin{proof}
Neem voor de richting van links naar rechts aan dat $Rxy$ geldt, en laat $z \in
\equivrep{x}{R}$. Per definitie geldt dan $Rxz$. Omdat $R$ een
equivalentierelatie is, geldt $Ryz$. (Volledig uitgeschreven: uit $Rxy$
en de symmetrie van~$R$ volgt $Ryx$, en uit dit omgekeerde verband,
$Rxz$ en de transitiviteit van~$R$ volgt~$Ryz$.) Dus $z \in
\equivrep{y}{R}$.
Hieruit volgt $\equivrep{x}{R} \subseteq \equivrep{y}{R}$. Op precies
dezelfde wijze volgt $\equivrep{y}{R} \subseteq \equivrep{x}{R}$.
Volgens extensionaliteit geldt dus $\equivrep{x}{R} =
\equivrep{y}{R}$.

Neem voor de richting van rechts naar links aan dat $\equivrep{x}{R} =
\equivrep{y}{R}$. Omdat $R$ reflexief is, geldt $Ryy$, en dus $y \in
\equivrep{y}{R}$. Dan volgt ook $y \in \equivrep{x}{R}$ uit de aanname
dat $\equivrep{x}{R} = \equivrep{y}{R}$. Dus geldt $Rxy$.
\end{proof}

\begin{ex}
Een mooi voorbeeld van equivalentierelaties komt uit het modulair rekenen.
Voor natuurlijke getallen $a$ en $b$ en voor $n \in \PosInt$ betekent
$a \equiv_n b$ dat
deling van $a$ door~$n$ dezelfde rest oplevert als deling van $b$ door~$n$.
(Iets symbolischer: $a \equiv_n b$ geldt precies wanneer er een $k \in
\Int$ is waarvoor $a - b = kn$.) $\equiv_n$ is voor iedere~$n$ een
equivalentierelatie. Bovendien zijn er precies $n$ verschillende
equivalentieklassen die door~$\equiv_n$ worden bepaald; dat wil zeggen
dat $\equivclass{\Nat}{\equiv_n}$
$n$ !!{element}s heeft. Dit zijn: de verzameling getallen die zonder rest
deelbaar zijn door $n$, namelijk $\equivrep{0}{\equiv_n}$; de verzameling
getallen die bij deling door $n$ rest~$1$ hebben, namelijk
$\equivrep{1}{\equiv_n}$; \ldots; en de verzameling getallen die bij deling
door~$n$ rest~$n-1$ hebben, namelijk~$\equivrep{n-1}{\equiv_n}$.
\end{ex}

\begin{prob}
Toon aan dat $\equiv_n$ voor iedere $n \in \PosInt$ een equivalentierelatie
is en dat $\equivclass{\Nat}{\equiv_n}$ precies $n$ elementen heeft.
\end{prob}
\end{document}
